\section{Political Economy Discussion and Conclusion}
\label{sec:discussion_conclusion}

\subsection{Competing Paradigms in Economic Thought: Evaluating Research Programmes on the Unidad Popular}
\label{sec:competing_paradigms}

\paragraph{Rival Worldviews on the Unidad Popular.}
Theoretical advance in economic history rarely proceeds through neutral empirical accumulation. For more than half a century, mainstream economics has interpreted the breakdown of Chile's Unidad Popular (UP, 1970--1973) through the framework of ``macroeconomic populism'' \citep{DornbuschEdwards1990,Edwards2024}. Under market-clearing assumptions and monetarist quantity theory, this literature attributes the 1972--1973 inflation to uncoordinated fiscal-monetary expansion and discretionary currency emission. Yet this orthodox consensus analyzes peripheral capitalism as a closed monetary system, abstracting from international currency hierarchies, external financial subordination, and balance-of-payments constraints. Evaluated through \citeauthor{Lakatos1978}'s (\citeyear{Lakatos1978}) methodology of scientific research programmes, the orthodox interpretation exhibits persistent anomalies, while dependency theory \citep{Kvangraven2021} provides a progressive framework capable of explaining the high-frequency empirical record.

\paragraph{The Monetarist Paradigm and Its Explanatory Anomalies.}
A research programme degenerates when its theoretical core produces no novel predictions and survives primarily by introducing ad hoc adjustments to accommodate anomalies \citep{Lakatos1978}. The orthodox monetarist programme rests on three foundational axioms: the exogenous determination of the money supply, direct quantity-theoretic pass-through from money to prices, and long-run monetary neutrality under full-capacity market clearing. Confronted with the historical record of 1971---when base money expanded by over 100\% while inflation fell from 35\% to 22\% and manufacturing output expanded by 12.1\%---orthodoxy relies on narrative rationalizations. It dismisses the cancellation of foreign commercial credit lines, the U.S. Eximbank embargo, and falling copper export prices as secondary disturbances \citep{Edwards2024}. In addition, its reliance on market-clearing assumptions leads it to diagnose retail shortages solely as ``repressed inflation'' caused by price controls, ignoring organized employer distribution lockouts and the fact that manufacturing capacity utilization reached 98.4\% alongside contracting gross fixed investment in machinery ($-2.3\%$).

\paragraph{Dependency Theory as an Open-Economy Framework.}
Dependency theory grounds macroeconomic dynamics in structural asymmetries: external financial subordination, specialized primary export structures, and an unequal international currency hierarchy \citep{Prebisch1950,Marini1973,DosSantos1970,Vasudevan2009,Basu2022}. This paper operationalizes this tradition by formalizing central bank balance-sheet solvency as an endogenous state variable ($S_t \equiv F_t / H_t$). This analytical framework accounts for empirical patterns that closed-economy orthodoxy cannot explain: the state-dependent breakdown of monetary insulation, the empirical dominance of reverse accommodation from prices to base money, and the dissolution of the external price buffer once foreign exchange reserves are exhausted.

\subsection{Empirical Causality and Dynamic Transmission: Precedence and Non-Linear Multipliers}
\label{sec:causality_transmission}

\paragraph{Directional Precedence and Endogenous Accommodation.}
The econometric findings in Section~\ref{sec:historical_empirical_results} provide a direct empirical test between these competing frameworks. Standard monetarist accounts treat central bank emission as an autonomous policy impulse that initiates price increases. The Granger causality tests within stationary vector autoregressions in Section~\ref{sec:granger_tournament} reject this assumption. Temporal precedence exhibits bidirectional predictive interaction with dominant reverse accommodation running from consumer price inflation to base money growth ($\pi \to g^H$). The Granger $F$-statistics for reverse accommodation are four times larger than those for forward money growth and persist across all tested lag horizons ($k=1,\dots,4$), while monthly solvency ratio growth systematically predicts base money expansion ($g_{\text{SolvR\_H}} \to g^H$). In parallel, the constraint-based PCMCI causal discovery algorithm in Appendix~\ref{app:pcmci_ee1} retains directed lagged conditional links connecting gross reserve depletion to subsequent base money expansion and relative price distortions. The Central Bank functioned primarily as an endogenous accommodation valve, issuing credit to sustain payrolls and prevent production halts as external trade credit froze.

\paragraph{Dynamic Multipliers and Reverse Accommodation Dominance.}
The non-linear Threshold Vector Autoregression (TVAR) and Generalized Impulse Response Function (GIRF) estimates in Section~\ref{sec:tvar_girf_multipliers} provide quantitative evidence against money exogeneity. Across both regimes, reverse defensive accommodation ($\pi_m \to g^H$) exceeds forward monetary transmission ($g^H \to \pi_m$) in persistence and cumulative response. In the adverse reserve-coverage regime ($\Delta s_{t-1} \le -4.615\%$), a one-standard-deviation base money shock produces a modest price response: domestic inflation rises by 0.97 percentage points at month 1 and yields a cumulative 12-month response of 3.55 percentage points (point estimate; 3.90 pp bootstrap median). In contrast, a one-standard-deviation price shock triggers an immediate, sustained monetary expansion: base money expands by 1.72 percentage points at month 1 and remains statistically significant for 10 consecutive months ($h=1$ to $10$), generating a cumulative 12-month emission of 6.54 percentage points (point estimate; 7.00 pp bootstrap median). Dynamic multiplier difference tests across regimes ($\Delta(h) \equiv \text{GIRF}^{\text{Insolv}}(h) - \text{GIRF}^{\text{Normal}}(h)$) establish that reverse accommodation escalates significantly under reserve distress, with cross-regime divergence ($\Delta_{\pi_m \to g^H}$) statistically significant across 10 consecutive months ($h=1$ to $10$). When enterprise costs escalated amidst supply disruptions, the monetary authority accommodated liquidity demand to prevent widespread bankruptcy.

\paragraph{State-Dependent Transmission: External Buffers and Reserve Deterioration.}
The TVAR estimates demonstrate that international price transmission is state-dependent on central bank reserve backing. In the non-adverse reserve regime, foreign reserves provide an effective external buffer: world price shocks ($g^{\text{gold}} \to \pi_m$) exert no statistically significant effect on domestic consumer prices over a 12-month horizon (Panel~A of Figure~\ref{fig:tvar_girf_gold_shield}). The central bank absorbs import cost pressures through foreign exchange liquidity. Once reserve depletion crosses the estimated structural threshold in the growth gap ($\Delta s_{t-1} \le -4.615\%$), this external buffer dissolves. In the adverse reserve-coverage regime, the identical world price shock generates an unbroken 10-month inflationary wave ($h=3$ to $12$), peaking at 0.98 percentage points at month 3 and producing 6.54 percentage points of cumulative inflation (point estimate; 6.93 pp bootstrap median). Point-wise multiplier difference tests ($\Delta_{g^{\text{gold}} \to \pi_m}$) confirm statistically significant cross-regime divergence across 10 consecutive months ($h=3$ to $12$). Endogenous regime-switching probabilities in Table~\ref{tab:app_d03_regime_switching} indicate that this adverse state exhibits high persistence ($\Delta P \approx 0$). Under depleted reserves, unbuffered international price shocks generate sustained domestic inflation, inducing defensive monetization to sustain production.

\subsection{The International Currency Hierarchy, Hegemonic Contingency, and Central Bank Agency}
\label{sec:currency_hierarchy_agency}

\paragraph{The Peripheral Paradox: Domestic Credit Tokens versus International Reserve Money.}
These dynamic responses reflect what I define as the \emph{Peripheral Paradox} (Section~\ref{sec:money_endogeneity_paradox}). The central bank balance sheet in dependent economies embodies two distinct monetary forms. On the liability side, domestic base money ($H_t$) consists of sovereign credit tokens issued under national authority \citep{Basu2022}. On the asset side, international reserves ($F_t$) represent claims on world money that the peripheral state cannot emit \citep{Vasudevan2009,AlamiEtAl2022}. When foreign commercial credit lines are severed and export receipts fall, the domestic monetary circuit decouples from international settlement. The central bank confronts an institutional dilemma: either it withholds liquidity from enterprises facing import price explosions, triggering immediate insolvency and mass layoffs; or it acts as lender of last resort, expanding credit that depreciates the currency in the absence of foreign exchange backing.

\paragraph{Hegemonic Contingency, Constrained Agency, and the Bretton Woods Environment.}
The Chilean transition unfolded against the breakdown of the post-war international monetary order. As \citet{Zoeller2019} demonstrates, President Nixon's decision to close the gold window on August 15, 1971, was a tactical contingency undertaken to manage U.S. balance-of-payments pressures, which exported stagflationary instability to the periphery \citep{Vernengo2021}. Simultaneously, Chile faced a targeted bilateral financial blockade: foreign commercial banks canceled over \$220~million in short-term credit lines, while the U.S. Eximbank and multilateral institutions halted disbursements \citep{DeVylder1976,PetrasMorley1975}. In this hostile context, the Central Bank of Chile exercised constrained institutional agency. Documented in its \textit{Memorias Anuales} and board records, the BCCh chose to accommodate enterprise credit demand and finance essential food imports, consciously defending employment and production rather than imposing a catastrophic deflationary adjustment on the working class.

\paragraph{Institutional Limits of Currency Sovereignty: A Clarification on Neo-Chartalism.}
This historical experience clarifies the institutional limits of currency sovereignty in the Global South. While neo-chartalist frameworks correctly emphasize that sovereign governments face no purely domestic budget constraints \citep{AlamiEtAl2022}, formal legal authority over domestic fiat money does not confer command over world money. In an asymmetric currency hierarchy, a peripheral government cannot overcome foreign exchange bottlenecks by issuing domestic credit. Without convertible reserves to purchase non-substitutable capital goods, spare parts, and food imports, domestic liquidity injections lead to currency depreciation and imported cost-push inflation. Progressive macroeconomic policy in the periphery remains fundamentally bounded by the external balance-of-payments constraint.

\subsection{Synthesis, Methodological Scope, and Concluding Remarks}
\label{sec:synthesis_horizons}

\paragraph{Evidentiary Limits and Econometric Caveats.}
A rigorous appraisal must acknowledge the evidentiary limits of macroeconomic time-series econometrics. The empirical methods deployed here---stationary vector-autoregressive Granger causality tests, PCMCI causal discovery, and non-linear threshold vector autoregressions---establish directional precedence and state-dependent conditional dynamics. They identify robust empirical patterns across historical regimes, but they do not constitute micro-level experimental counterfactuals. Macro aggregates remain subject to unobserved common shocks and structural breaks. Furthermore, the 1960--1980 sample pools the Unidad Popular with the Christian Democratic and early military dictatorship periods, and monthly data cannot directly observe nominal wage readjustments or strike activity. These econometric findings provide quantitative evidence consistent with balance-sheet constraints and endogenous accommodation, but they should be viewed as opening the empirical inquiry into peripheral monetary institutions.

\paragraph{Analytical Findings and Directions for Future Research.}
By grounding the macroeconomic history of the Unidad Popular in an open-economy, balance-sheet framework, this paper demonstrates the analytical utility of dependency theory as a research programme \citep{Kvangraven2021}. The analysis accounts for empirical anomalies that disrupt orthodox narratives, replacing closed-economy quantity-theoretic assumptions with an empirical model of reserve solvency. This framework points toward several avenues for future research. One important frontier involves examining the meso-scale institutional mechanisms of credit allocation, tracing how liquidity flowed between the central bank, nationalized banks, and state enterprises during balance-of-payments crises. Another promising avenue is comparative: exploring whether the solvency threshold and reserve buffering dynamics identified in Chile operate across other peripheral economies confronting external debt and currency crises.

\paragraph{Institutional Implications: Foreign-Exchange Management and Structural Transformation.}
The historical experience of the Unidad Popular points beyond technical macroeconomic management. Foreign-exchange budgeting, capital controls, and the regulation of external financial flows can widen policy space and mitigate exposure to international liquidity cycles \citep{FfrenchDavis2018}. Yet macroprudential management within the balance-of-payments constraint does not by itself dissolve the structural relations that reproduce that constraint. A peripheral state may command the domestic means of payment while remaining unable to produce the world money required to reproduce an import-dependent productive structure. Foreign exchange is therefore not merely an instrument of monetary administration; it is a material condition of accumulation and a focus of structural conflict over the control of export revenues, external credit, and imported capital inputs. Macroprudential autonomy can cushion the shocks of dependent accumulation, but the fundamental challenge of peripheral transformation remains altering the structural relations through which external vulnerability is reproduced.
